\section{Methodology}\label{sec:metodologia}

This section details the analysis and optimization framework developed for the parametric design of roundwood timber bridges, specifically of the girders and deck planks.

The framework was implemented in \textit{Python} and made freely available through the \textit{ReliaBridge} platform \parencite{ReliaBridge2026}.
Metaheuristic multi-objective optimization is carried out with the genetic algorithms of the \texttt{pymoo} library \parencite{blank2020}, while uncertainty quantification and stochastic global sensitivity analysis are performed with the \texttt{UQpy} library \parencite{Olivier2020}.

For each processed solution, the platform provides the complete calculation report, indicating the internal forces, the displacements and the outcome of each ultimate and serviceability code check.

The evaluation flow of each candidate solution follows the structure already presented in Algorithm~\ref{alg:nsga2} (Section~\ref{sec:nsga2}). For each individual $\mathbf{x}$ generated by NSGA-II, the diameter is perturbed over a grid of $N_c$ points and each realization is submitted to the \texttt{EvaluateModel} function. The objective functions remain computed at the nominal point $\mathbf{x}$, without any aggregation over the grid, whereas the constraints are aggregated by the worst-case operator (maximum, component-wise) over the $N_c$ perturbed realizations, as described in Section~\ref{sec:robustez}.
Algorithm~\ref{alg:avaliarmodelo} details the internal procedure of this function for a single design vector, nominal or perturbed.

\begin{algorithm}[!ht]
\caption{EvaluateModel($\tilde{\mathbf{x}}$)}
\label{alg:avaliarmodelo}
\begin{algorithmic}[1]
\Require Design vector $\tilde{\mathbf{x}} = (d,\ b_w,\ h,\ esp_{long},\ esp_{tab})^T$
\Ensure Objective vector $\mathbf{f} = (f_1, f_2)$ and constraint vector $\mathbf{g} = (g_1,\ldots,g_6)$
\State Run the geometric accommodation routine, determine $n_{long}$ and $n_{tab}$ and evaluate the spacing constraints $g_5$ and $g_6$
\State Define the permanent and variable actions (design vehicle and crowd load) acting on the deck
\State Distribute the deck reaction onto the girders (couples the loading of both members into a single girder loading)
\State \textbf{Girder:} compute the internal forces ($M_{sd}$, $Q_{sd}$) and deflections, and evaluate $g_1$ (bending), $g_2$ (shear) and $g_3$ (deflection, with creep)
\State \textbf{Deck:} compute the bending moment under the design vehicle wheel and evaluate $g_4$ (bending)
\State Compute $f_1$ (total timber volume, Eq.~\eqref{eq:f1}) and $f_2$ (deflection utilization, Eq.~\eqref{eq:f2})
\State \Return $\mathbf{f}, \mathbf{g}$
\end{algorithmic}
\end{algorithm}

Step 3 highlights the coupling between deck and girder that characterizes the formulation. The deck reaction, arising from the permanent loads and the design vehicle, is distributed directly onto the girders at each evaluation, so that both members are designed simultaneously rather than in independent steps.

\subsection{Test model definition and material properties}\label{sec:modelo_teste}

To test the tool and demonstrate the direct application of the implemented algorithm, a single-lane rural road bridge designed to carry the TB-240 design vehicle traffic specified by \textcite{nbr7188} was adopted as the case study.
This configuration is representative of rural accesses and local roads, where the demand for timber designs is high.

The mechanical modelling of the structural members strictly follows the Brazilian code provisions.
The material characteristics and the properties adopted in the simulations were established in accordance with \textcite{NBR7190-1:2022}.
The elastic properties and the base strength classes for native forest species (D20 to D60), derived from tests on defect-free specimens, are listed in Table~\ref{tab:classes_madeira}.

It should be noted that, according to item 6.3.4 of \textcite{NBR7190-1:2022}, when designing with the classes of Table 2 of the code, $f_{m,d} = f_{c0,d}$ is assumed, that is, the design bending strength is taken equal to the design compressive strength parallel to the grain.
Since the modification factor and the partial factor for normal stresses ($\gamma_{wf}$) are the same for both checks, this equivalence translates, at the characteristic level, into $f_{m,k} = f_{c0,k}$.
Hence, the bending strengths used by the platform in the limit state constraints are the tabulated $f_{c0,k}$ values of each class.

\begin{table}[!ht]
\caption{Strength classes of native forest species and computed bending strength.\label{tab:classes_madeira}}
\begin{threeparttable}
\footnotesize
\setlength{\tabcolsep}{3pt}
\begin{tabular*}{\columnwidth}{@{\extracolsep{\fill}}lccccc@{\extracolsep{\fill}}}
\toprule
Class & $f_{c0,k}$ (MPa) & $f_{m,k}$ (MPa) & $f_{v0,k}$ (MPa) & $E_{c0,med}$ (MPa) & $\rho_{12\%}$ (kg/m\textsuperscript{3}) \\
\midrule
D20 & 20 & 20.0 & 4 & 10\,000 & 500 \\
D30 & 30 & 30.0 & 5 & 12\,000 & 625 \\
D40 & 40 & 40.0 & 6 & 14\,500 & 750 \\
D50 & 50 & 50.0 & 7 & 16\,500 & 850 \\
D60 & 60 & 60.0 & 8 & 19\,500 & 1\,000 \\
\bottomrule
\end{tabular*}
\end{threeparttable}
\end{table}
\subsection{Structural system and parametric model}\label{sec:sistema}

The structural system considered in the computational model consists of a deck made of longitudinal rectangular timber planks, continuously supported on longitudinal circular roundwood beams (girders), regularly distributed at a distance called the parametric spacing ($esp$).
The deck is subjected to permanent actions from the self-weight of its constituent materials and to variable crowd and axle actions associated with moving traffic.
The load transfer conveys the deck internal forces to the supporting girders, which are analysed as distributed loads along the span.
The conceptual representation of the system is shown in Fig.~\ref{fig:sistema}.

\begin{figure}[!ht]
\centering
\includegraphics[width=0.71\textwidth]{figuras/caminhao.png}
\caption{Conceptual representation of the structural system of the timber bridge.}
\label{fig:sistema}
\end{figure}

The model is mathematically parameterized. From a user input vector comprising the roadway width, mechanical properties, creep coefficients ($\varphi$) and code combination factors ($\psi_2$), the geometric properties of the cross-sections (area, second moment of area, section modulus and radii of gyration) are dynamically recomputed in the code for any generated individual.
This parametric flexibility couples the geometry to the formulation of the geometric constraints, enabling feasibility checks in automated computational routines.

\subsection{Constraint formulation and robust optimization}\label{sec:formulacao}

Within NSGA-II, the design is converted into a vector of mathematical variables $\mathbf{x} = (d,\allowbreak\ b_w,\allowbreak\ h,\allowbreak\ esp_{long},\allowbreak\ esp_{tab})^T$, containing the girder diameter ($d$), the deck plank dimensions ($b_w$, $h$), the spacing between girders ($esp_{long}$) and the spacing between deck planks ($esp_{tab}$).
From the spacings and the roadway width, the number of members of each family is determined internally by the geometric accommodation routine, which also gives rise to constraints $g_5$ and $g_6$.

Both objective functions are to be minimized.
The first, $f_1$, is the total timber volume of the superstructure, according to Eq.~\eqref{eq:f1}, where $n_{long}$ and $n_{tab}$ are the numbers of girders and deck planks actually accommodated, $A_{long}$ and $A_{tab}$ are the respective cross-sectional areas and $b_{w,\mathrm{pista}}$ is the roadway width.
The second, $f_2$, is the utilization of the girder serviceability deflection limit, defined in Eq.~\eqref{eq:f2} as the ratio between the total deflection, already including creep, and the code deflection limit $L/350$.
Minimizing $f_2$ is equivalent to maximizing serviceability performance, and the two functions are conflicting, since reducing the timber volume necessarily increases the relative deflection.

\begin{equation}
f_1(\mathbf{x}) = V_{\mathrm{total}} = n_{long}\,A_{long}\,L + n_{tab}\,A_{tab}\,b_{w,\mathrm{pista}}
\label{eq:f1}
\end{equation}

\begin{equation}
f_2(\mathbf{x}) = \frac{\delta_{\mathrm{tot}}}{L/350}
\label{eq:f2}
\end{equation}

The generated solutions must satisfy a strict set of six normalized constraint equations $g_j \leq 0$, namely girder bending stress ($g_1$), girder shear ($g_2$), girder deflection including creep ($g_3$), deck bending stress under the design vehicle wheels ($g_4$) and the code minimum and maximum spacing limits of the superstructure members ($g_5$ and $g_6$).
The code procedure for computing the geometric properties, internal forces and the bending, shear and deflection checks of the girder, as well as the bending check of the deck, is detailed in Section~\ref{sec:projeto_ponte}.

Owing to the high natural variability of roundwood diameter, the platform introduces robustness by evaluating $N_c$ realizations for each individual, perturbing only the girder diameter over a deterministic grid with deviation $\rho = 5\%$ (Section~\ref{sec:otimizacao}).

To measure the impact that the dimensional uncertainty of each of the five design variables, over its whole search domain, has on the six listed code constraints, the system integrates the computation of total-order Sobol sensitivity indices ($S_{T_i}$), evaluated by Monte Carlo simulation with the \texttt{UQpy} library \parencite{Olivier2020}.
This index numerically indicates to the designer which geometric characteristic of the bridge requires tighter quality control in cutting and field execution to prevent systemic failure of the deck.
